% !PRINTED 226
\weberpage{226}
%<<CONT
one therefore only has to ensure that the $\nu$ values of (3), which correspond to the $\nu$ rows, are distinct from one another, \gesperrt{so that (3) becomes the root of an (irreducible) transformation equation}\textsuperscript{1)}.

We assume that the function $\Phi(x)$ is either an \gesperrt{even} or an \gesperrt{odd} function of $x$, and accordingly make the following distinction.

If $\Phi(x)$ is an even function, then among the factors of the product (3) two at a time, namely

\weq{4}{\Phi(\mathrm{sn}\,h\,\Omega), \quad \Phi[\mathrm{sn}\,(n-h)\,\Omega]}

are equal to one another. If we therefore set

\weq{5}{\Pi(\Omega) = \prod_{1,\,\frac{n-1}{2}}^{h} \Phi(\mathrm{sn}\,h\,\Omega),}

then, if $m$ is any number coprime to $n$,

\weq{6}{\Pi(m\,\Omega) = \Pi(\Omega)}

[because among the $\frac{1}{2}(n-1)$ numbers $h\,m$ no two have a sum or a difference divisible by $n$, so that, apart from sign, the $\mathrm{sn}\,h\,\Omega$ are the same as the $\mathrm{sn}\,h\,m\,\Omega$]. Thus $\Pi(\Omega)$ is the root of a transformation equation. This class of transformation equations we call \gesperrt{modular equations}.

If $\Phi(x)$ is an odd function, then the two quantities (4) are opposite to one another and (6) is no longer generally correct; instead we have

\weq{7}{\Pi(m\,\Omega) = \pm\,\Pi(\Omega).}

Hence, at least in general, it is no longer $\Pi(\Omega)$, but only $\Pi(\Omega)^2$, that is a root of a transformation equation. This kind of transformation equations we call \gesperrt{multiplier equations}. In order to learn the cases in which $\Pi(\Omega)$ itself is a root of a transformation, the sign in (7) must be determined.

This is achieved by means of a theorem of number theory which is derived in Bd.\ I, \S\,145.

The theorem runs:
%CONT>>

\textsuperscript{1)}~Expressions like (3) are still roots of transformation equations even when $h$ takes only values coprime to $n$. Such a transformation equation % [sic] The German source prints the plural determiner with a singular noun.
has hitherto been little used. We shall become acquainted with an example further on.
